% А2.1 — Степень с целым показателем, свойства степеней
% Старая методичка: стр. 58; 65 (старые задания 6; 7)
% Источники: методичка «Степени» (определения, свойства 1–13, задача 3),
% шпаргалка «Степень». Стиль — стиль-изложения.md

\begin{tema}{А2.1}
\opr{Степенью} числа $a$ с натуральным показателем $n\geqslant2$ называется произведение
$n$ множителей, каждый из которых равен $a$:
\quad$a^n=\underbrace{a\cdot a\cdot\ldots\cdot a}_{n\ \text{раз}}$.\quad
Число $a$ — \emph{основание}, $n$ — \emph{показатель}. По определению $a^1=a$.

\opr{Нулевой и отрицательный показатель.} При $a\ne0$: \ $a^0=1$, \ $a^{-n}=\dfrac1{a^n}$.
Например, $2^{-3}=\frac18$, $\left(\frac23\right)^{-2}=\left(\frac32\right)^2=\frac94$.
Выражения $0^0$ и $0^{-n}$ не имеют смысла.

\opr{Знак степени отрицательного числа} зависит от чётности показателя:
$(-a)^n=a^n$, если $n$ чётное, и $(-a)^n=-a^n$, если $n$ нечётное.
Например, $(-2)^4=16$, $(-2)^3=-8$.

\begin{pfTablica}[bez shapki]{ll}
\pfPodshapka{Свойства степеней (далее $a\ne0$, $b\ne0$, $m$ и $n$ — целые)}\\
$a^m\cdot a^n=a^{m+n}$ & $\dfrac{a^m}{a^n}=a^{m-n}$\\
$(a^m)^n=a^{mn}$ & $1^n=1$\\
$(ab)^n=a^n b^n$ & $\left(\dfrac ab\right)^n=\dfrac{a^n}{b^n}$\\
\end{pfTablica}

Свойства работают в обе стороны. Главный приём в вычислениях —
\emph{привести всё к одному основанию}: $8=2^3$, $0{,}25=2^{-2}$, $\frac1{27}=3^{-3}$.

\begin{primer}
Найдите значение выражения $\dfrac{5^5\cdot121^{-3}}{11^{-7}\cdot125^2}$.
Решение: приведём к основаниям 5 и 11: $121=11^2$, $125=5^3$. Получим
$\dfrac{5^5\cdot11^{-6}}{11^{-7}\cdot5^6}=5^{5-6}\cdot11^{-6-(-7)}=5^{-1}\cdot11=\frac{11}5$.
Ответ: 2,2.
\end{primer}

\begin{oshibka}
$-2^4=-16$, а $(-2)^4=16$: без скобок в степень возводится только 2.
Показатели складываются только при одинаковых основаниях: $2^3\cdot3^2\ne6^5$.
\end{oshibka}
\end{tema}
